Une bouilloire électrique branchée sous \qty{120}{V} a une puissance de \qty{\v{P}}{W}. Elle contient \qty{\v{m}}{L} d'eau à \qty{\v{Ti}}{\celsius} qu'on veut amener à \qty{100}{\celsius} (pour l'eau, $c = \qty{4186}{J/kg.\celsius}$ et \qty{1}{L} a une masse de \qty{1}{kg}). On suppose que toute l'énergie électrique sert à chauffer l'eau.

\begin{enumerate}[label=\alph*)]
    \item Déterminez la résistance de l'élément chauffant et le courant qui le traverse.
    \item Déterminez l'énergie nécessaire pour amener l'eau à ébullition.
    \item Combien de temps faut-il pour amener l'eau à ébullition ?
\end{enumerate}
